\documentclass[]{article}

\usepackage{amsmath}
\usepackage{multirow}
\usepackage{natbib}
\usepackage{graphicx} 
\usepackage[margin=1in]{geometry}

\begin{document}

% Auto-generated figures file
% \usepackage{graphicx}

\begin{figure}[h!]
    \centering
    \includegraphics[width=0.99\textwidth]{../input_files/plots/iter2_conditional_Rv_1779134732.png}
    \caption{Normalised Lagrangian velocity autocorrelation function, R\_v(τ), for the full tracer ensemble (black), the strain-dominated subpopulation (>70\% time in strain, red), and the vortex-dominated subpopulation (<30\% time in strain, blue). The slow decay indicates a long velocity correlation time, which drives the initial quasi-ballistic transport. The curves for the strain-dominated and vortex-dominated populations are nearly identical, refuting the hypothesis of distinct transport mechanisms and showing that long-term velocity correlations are homogeneous throughout the flow.}
    \label{fig:image0}
\end{figure}

\begin{figure}[h!]
    \centering
    \includegraphics[width=0.99\textwidth]{../input_files/plots/iter2_msd_Ht_1779134732.png}
    \caption{Mean Squared Displacement (MSD) as a function of lag time τ (left) and the temporal evolution of the running Hurst exponent H(t) (right) for the full tracer ensemble, strain-dominated (>70\% time in strain), and vortex-dominated (<30\% time in strain) sub-populations, compared to a phase-randomised surrogate. The analysis reveals that all tracer populations exhibit nearly identical transport properties, with H(t) decaying from a near-ballistic regime (H ≈ 1) to a near-diffusive regime (H ≈ 0.56) at late times. This demonstrates a pre-asymptotic crossover and refutes the hypothesis of a distinct, ballistically transported strain-dominated population. In contrast, the surrogate maintains a near-ballistic exponent (H ≈ 0.98), indicating that strong restorative temporal correlations in the original velocity field suppress long-range transport.}
    \label{fig:image1}
\end{figure}

\begin{figure}[h!]
    \centering
    \includegraphics[width=0.99\textwidth]{../input_files/plots/iter2_strain_fraction_H_1779134732.png}
    \caption{Refutation of the strain-dominated highway hypothesis. (Left) Bimodal distribution of the 8000 tracers based on the percentage of time spent in strain-dominated regions, defining vortex-dominated ($<30\%$) and strain-dominated ($>70\%$) sub-populations. (Right) The late-time Hurst exponent H is nearly identical for the full ensemble (H = 0.565), the strain-dominated sub-population (H\_strain = 0.557), and the vortex-dominated sub-population (H\_vortex = 0.574). All values are close to the normal diffusion limit (H = 0.5), demonstrating that long-time transport scaling is independent of the tracer's history and refuting the existence of a ballistic transport mechanism.}
    \label{fig:image2}
\end{figure}



\end{document}
                